%----------------------------------------------------------------------------
\chapter{Kutatási eredmények és kiértékelés}
%----------------------------------------------------------------------------

A dolgozat ezen részében egy rövid távú felmérés segítségével értékeljük ki a megvalósított rendszert. A vizsgálat célja az, hogy megértsük, milyen mértkben és módon támogatja az AI alapú alkalmazás használata a tanulási folyamatot, témaspecifikus szókincsfejlődést, beszédgyakorlést, valamint a tanulók motivációját egy hagyományos módszerhez képest.

A kuttás alatt 16 tanulót mértünk, 8 tanuló használta a hagyományos módszert, míg 8 tanuló az AI alapú alkalmazást. Az értékelés során előteszt, utóteszt, beszédkészség mérés került feldolgozásra. Az tesztek eredményeinek összehasonlítására leíró statisztikai elemzést végeztem, amely magában foglalja az átlagok, szórások és a különbségek vizsgálatát a két csoport között. Emellett a csoportokon belüli fejlődést párosított t-próbával értékeltem, hogy megállapítsam, van-e szignifikáns javulás a tanulók teljesítményében a tesztek között, illetve a két csoport közötti fejlődés különbséget független mintás t-próbával vizsgáltam.

Az mérés mivel rövid távú, nem ad nekünk teljes képet a használhatóságról és a hosszú távú hatásokról, hanem inkább az a célja hogy megvizsgáljuk vele, hogy működőképes-e az oktatási környezetben, és milyen tanulói visszajelzéseket kapunk a használat során.

\section{A vizsgálat lebonyolítása}

A vizsgálatban 16 tanuló vett részt, akik egy magántanár diákjai. 8 hagyományos oktatási módszerrel tanulta ugyanazt a témát, míg a másik 8 tanuló az alkalmazással tanulta.

A tanulási folyamat minden diáknak 40 percet vett igénybe. A tanóra témája a technológiahasználat volt, amely során a tanulók az online biztonságróé, kártékony szoftverekről és a digitális eszközök használatáról tanultak.

A kontrollcsoport tanulási folyamata egy hagyományosabb módszerrel zajlott, amely során a tanár prezentációt tartott a témáról, majd vocabulary feladatokkal és beszélgetéssel illetve brainstorminggal folytatták. A kísérleti csoport ezzel ellentétben az alkalmazásban olvasta el a lecke szövegét, vagy felolvastathatta az AI asszisztenssel is, illetve a szavakat is felolvastathatta vele, hogyha nem tudta a kiejtést. Ebben a szövegben kijelölhette az ismeretlen szavakat vagy kifejezéseket, amjd AI asszisztnssel beszélgetett a témában. A beszélgetés során az alkalmazásban tárolódtak a hibák, amelyeket beszélgetés után flashcardok segítségével vagy javításos feladatokkal gyakorolhattak.

A mérés mindkét csoport eseténben az óra elötti előtesztből, óra utáni utótesztből és a beszédkészség méréséből állt, miután kérdőívvel gyűjtöttük be a visszajelzéseket. Az előteszt segített megadni a kiindulási pontjukat a diákoknak, az utóteszt pedig a tanulás hatékonyságát mérte. Ezzel mértük a tanulók fejlősését, élményét, motivációját, és használhatóságát az alkalmazásnak.

\section{A csoportok előteszt eredményeinek összehasonlítása}

Az előteszt célja annak felmérése, hogy a kontrollcsoport és a kísérleti csoport tanulói milyen kiindulási szinttel rendelkeznek a tanulás megkezdése előtt. Ez azért nagyon fontos, mivel két tanulási módszer összehasonlításához elengedhetetlen szempont az, hogy mindkét csoport hasonló kiindulási ponttal rendelkezzen nagyjából.

Az előteszten 20 pontot lehetett elérni, ez a tanóra témájához kapcsolódó szókincset, fogalomértést vizsgálta. A feladatok között szerepelt a szó-jelentés párosítás, fill in the gap, valamint kérdés-válasz típusú feladatok. A két csoport eredményeit leíró statisztikai adatotait a~\ref{tab:eloteszt-leiro-stat} táblázat mutatja be.

\begin{table}[h!]
    \centering
    \begin{tabular}{lccc}
        \hline
        \textbf{Csoport} & \textbf{Létszám} & \textbf{Előteszt átlaga} & \textbf{Szórás} \\
        \hline
        Kontrollcsoport & 8 & 5.75 & 1.19 \\
        Kísérleti csoport & 8 & 5.62 & 2.17 \\
        \hline
    \end{tabular}
    \caption{Az előteszt eredményeinek leíró statisztikái csoportonként}
    \label{tab:eloteszt-leiro-stat}
\end{table}

A táblázat adatai alapján látható, hogy a két csoport előteszt eredményei közel hasonló eredményekkel rendelkeznek, a kontrollcsoport átlaga 5.75, míg a kísérleti csoport átlaga 5.62. Ez arra utal, hogy a tanulók hasonló szinttel rendelkeznek, így a két csoport összehasonlítható kiindulási tudásszinttel rendelkezik.

\begin{table}[h!]
    \centering
    \renewcommand{\arraystretch}{1.2}
    \begin{tabular}{p{4cm}cccc}
        \hline
        \textbf{Csoport} & \textbf{t érték} & \textbf{Szabadságfok} & \textbf{p-érték} & \textbf{Értelmezés} \\
        \hline
        Kontrollcsoport előteszt vs. kísérleti csoport előteszt & 0.13 & 14 & 0.896 & Nincs jelentős különbség \\
        \hline
    \end{tabular}
    \caption{Az előteszt t-próba eredményei a két csoport között}
    \label{tab:eloteszt-t-proba}
\end{table}

A \ref{tab:eloteszt-t-proba} táblázatban látható, hogy a függetelen mintás t-próba eredményei alapján nincs szignifikáns különbség a két csoport előteszt eredményei között. Az előteszt eredmények alapján tehát nem mutatható ki statisztikailag szignifikáns különbség a két csoport között. t(14)=0.13, p=0.896. Ez arra utal, hogy a kéy csoport hasonló kiindulási szinttel rendelkezik a vizsgálat kezdetén, így a összehasonlíthatónak tekinthetők.

\section{Előteszt és utóteszt eredményei csoporton belül}

A következő lépésben a csoportokon belüli fejlődést vizsgáltam meg. Mindkét csoport esetében összehasonlítottam az előteszt és utóteszt eredményeit. Az utóteszt szintén 20 pontos, előteszthez hasonló feladatokkal rendelkezik, de nem ugyanazokkal. A két mérés különbségei alapján értékeltem a tanulók fejlődését a tanóra során.

A két csoport csoporton belüli fejlődését a~\ref{tab:eloteszt-utoteszt-fejlodes} táblázatban foglalom össze, amely a párosított t-próba eredményeit mutatja be.

\begin{table}[h!]
\centering
\begin{tabular}{lccc}
\hline
\textbf{Csoport} & \textbf{Előteszt átlaga} & \textbf{Utóteszt átlaga} & \textbf{Átlagos fejlődés} \\
\hline
Kontrollcsoport & 5.75 & 13.25 & 7.50 \\
Kísérleti csoport & 5.63 & 14.25 & 8.63 \\
\hline
\end{tabular}
\caption{Az előteszt és utóteszt eredményeinek átlagai és az átlagos fejlődés csoportonként}
\label{tab:eloteszt-utoteszt-fejlodes}
\end{table}

A fent látható táblázat alapján megfigyelhető, hogy mindkét csoport jelentős fejlődést mutatott az előteszthez képest. A kontrollcsoport eredményei 5.75-ről 13.25-re javultak, míg a kísérleti csoport eredményei 5.63-ről 14.25-re javultak. Így megfigyelhetjük, hogy mindkét módszer hasonlóan teljesített, a kísérleti csoport minimálisan magasabb szintet ért el, de a különbség nem jelentős. Ez azt jelenti, hogy az alkalmazás használata is támogatta a tanulók fejlődését.

A csoporton belüli fejlődés statisztikai vizsgálatára párosított mintás t-próbát végeztem el, amelynek eredményeit a~\ref{tab:eloteszt-utoteszt-t-proba} táblázat mutatja be. Ezt azért végeztem, hogy megállapítsam, van-e szignifikáns javulás a tanulók teljesítményében a tesztek között.

\begin{table}[h!]
\centering
\begin{tabular}{lcccc}
\hline
\textbf{Csoport} & \textbf{Összehasonlítás} & \textbf{t érték} & \textbf{df} & \textbf{p érték} \\
\hline
Kontrollcsoport & előteszt--utóteszt & 8,66 & 7 & $p < 0{,}001$ \\
Kísérleti csoport & előteszt--utóteszt & 20,54 & 7 & $p < 0{,}001$ \\
\hline
\end{tabular}
\caption{Az előteszt és utóteszt eredményeinek összehasonlítása párosított mintás t-próbával}
\label{tab:eloteszt-utoteszt-t-proba}
\end{table}

A táblázat adatai alapján mindkét csoport esetében szignifikáns javulás figyelhető meg az előteszt és utóteszt eredményei között. Ez azt jelenti, hogy a tanulók a tanulási szakasz után jelentős fejlődést mutattak, függetlenül attól, hogy melyik módszert használták.
A kontrollcsoport fejlődése azt jelzi, hogy a hagyományos prezentációs, vocabulary és interaktív brainstorming módszer is hatékony a gyerek szókincsének bővítésére és fejlődédésére a témában. A kísérleti csoport esetében pedig azt mutatja, hogy az alkalmazás által tanult lecke feldolgozása, ismeretlen szavak és kifejezések gyakorlása, AI beszélgetések és hibák újragyakorlása szintén támogatta a tanulást.

Összességében elmondható, hogy mindkét csoport jelentős fejlődést mutatott, és hogy nem mutatkozott kifejezetten nagy különbség a két módszer között. Ez arra utal, hogy az alkalmazás a kísérletünk szerint hasonló tanulási eredményeket tud nyújtani, miközben a tanulók visszajelzései alapján használható és motiváló.

\newpage

\section{A két csoport fejlődésének összehasonlítása}

A csoportokon belüli fejlődés vizsgálata után elengedhetetlen lépés az, hogy megvizsgáljuk a kontrollcsoport és a kísérleti csoport fejlődése közötti különbséget. Itt minden tanuló esetében kiszámításra került a fejlődési érték, amely az utóteszt és előteszt különbsége.

Az fejlődési értékek összesítését a~\ref{tab:fejlodes-osszesites} táblázat mutatuja be.

\begin{table}[h!]
\centering
\begin{tabular}{lccccc}
\hline
\textbf{Csoport} & \textbf{N} & \textbf{Átlag} & \textbf{Szórás} & \textbf{Min.} & \textbf{Max.} \\
\hline
Kontrollcsoport & 8 & 7,50 & 2,45 & 4 & 11 \\
Kísérleti csoport & 8 & 8,63 & 1,19 & 7 & 10 \\
\hline
\end{tabular}
\caption{Az átlagos fejlődés leíró statisztikái csoportonként}
\label{tab:fejlodes-osszesites}
\end{table}

A a~\ref{tab:fejlodes-osszesites} táblázat alapján a kísérleti csoport átlagos fejlődése valamivel magasabb volt mint a kontrolllcsoportté. 

Végül a két csoport fejlődési értékeinek összehasonlítására független mintás t-próbát végeztem el, amelynek eredményeit a~\ref{tab:fejlodes-t-proba} táblázat mutatja be.


\begin{table}[h!]
\centering
\begin{tabular}{p{6cm}cccc}
\hline
\textbf{Összehasonlítás} & \textbf{t érték} & \textbf{df} & \textbf{p érték} & \textbf{Értelmezés} \\
\hline
Kontrollcsoport fejlődése vs. kísérleti csoport fejlődése & -1,17 & 14 & 0,262 & nincs szignifikáns különbség \\
\hline
\end{tabular}
\caption{A kontrollcsoport és a kísérleti csoport fejlődésének összehasonlítása független mintás t-próbával}
\label{tab:fejlodes-fuggetlen-tproba}
\end{table}

A fent látható táblázat szerint (\ref{tab:fejlodes-fuggetlen-tproba}) a független mintás t-próba alapján a kontrollcsoport és kísérleti csoport fejlődése között nincs szignifikáns különbség, mivel a p érték nagyobb, mint 0.05. Ez azt jelenti, hogy a két csoport fejlődése hasonló mértékű volt, és nem mutatott szignifikáns különbséget a tanulási módszerek között. Bár a kísérleti csoport átlagos fejlődése valamivel magasabb volt, ez a különbség ilyen rövid távú vizsgálatban nem tekinthető statisztikailag jelentősnek.

Így ez azt jelenti, hogy az alkalmazás használata rövid távon hasonló tanuási eredményeket nyújt, mint a hagyományos, viszont ez nem tekinthető garantáltnak hosszú távra is. 

\newpage

\section{A kérdőíves visszajelzések összefoglalása}

\begin{table}[h!]
\centering
\begin{tabular}{p{5cm} p{9cm}}
\hline
\textbf{Vizsgált terület} & \textbf{Eredmény} \\
\hline
Használhatóság & az alkalmazás használata könnyűnek bizonyult \\
Tanulási élmény & a tanulók érdekesnek találták az AI-alapú gyakorlást \\
Vocabulary funkció & hasznos volt az ismeretlen szavak kijelölése és mentése \\
AI-beszélgetés & segítette a témához kapcsolódó szókincs aktív használatát \\
Hibajavítás & a tanulók hasznosnak találták a beszélgetés utáni visszajelzéseket \\
Motiváció & az alkalmazás pozitív tanulási élményt biztosított \\
\hline
\end{tabular}
\caption{A kérdőíves visszajelzések összefoglalása}
\label{tab:visszajelzesek}
\end{table}

Az alábbi~\ref{tab:visszajelzesek} táblázatban összefoglalom a kérdőíves visszajelzéseket, amelyeket a tanulóktól gyűjtöttem be a vizsgálat során. A kérdőív célja az volt, hogy megértsük a tanulók élményeit, motivációját és az alkalmazás használhatóságát. 

A kérdőíves visszajelzések alapján a tanulók mind pozitív élményről számoltak be az alkalmazás használata során. A tanulós könnyen tudták használni az alkalmazást, pozitív elemként említették az AI beszélgetéseket, illetve az ismeretlen szavak kijelölését és a hibák javítását is. A visszajelzések alapján az alkalmazás egyik legfontosabb előnye nem feltétlenül abban rejlik, hogy azonnal jobb teszteredményeket biztosítson, hanem az, hogy moticálóbb és interaktívabb, modernebb tanulási környezetet nyújtott.

\section{Hipotézisek értékelése}

A kutatási tervben megfogalmazódott néhány hipotézis, amelyeket a vizsgált eredmények alapján ebben a fejezetben értékelek ki. A hipotézisek értékelése hozzájárul ahhoz, hogy összefoglaljuk azt, milyen mértékben támasztotta alá az alkalmazással kapcsolatos előzetes feltételezéseket a rövid távú vizsgálat eredménye.

\begin{table}[h!]
\centering
\begin{tabular}{p{6cm} p{3cm} p{6cm}}
\hline
\textbf{Hipotézis} & \textbf{Eredmény} & \textbf{Értelmezés} \\
\hline
H1: Az AI-alapú alkalmazást használó tanulók témaspecifikus szókincse fejlődik az előteszt és az utóteszt eredményei alapján. 
& igazolódott 
& a kísérleti csoport előteszt-átlaga 5,63 pontról 14,25 pontra emelkedett, a fejlődés szignifikáns volt \\
\hline
H2: Az AI-alapú alkalmazást használó csoport fejlődése nagyobb mértékű lesz, mint a hagyományos módszerrel tanuló kontrollcsoporté. 
& nem igazolódott egyértelműen 
& az AI-csoport átlagos fejlődése magasabb volt, de a különbség nem volt statisztikailag szignifikáns \\
\hline
H3: Az AI-alapú alkalmazást használó tanulók beszédkészsége javul az adott témáról folytatott irányított beszélgetés során. 
& részben igazolódott 
& az AI-beszélgetés támogatta a témához kapcsolódó szókincs aktív használatát, de ehhez nem készült külön statisztikai próba \\
\hline
H4: A tanulók pozitívan értékelik az AI-alapú alkalmazás használhatóságát, érthetőségét és motiváló hatását. 
& igazolódott 
& a tanulói visszajelzések alapján az alkalmazás könnyen használhatónak, érthetőnek és motiválónak bizonyult \\
\hline
\end{tabular}
\caption{A kutatási hipotézisek értékelése}
\label{tab:hipotezisek_ertekelese}
\end{table}

A~\ref{tab:hipotezisek_ertekelese} táblázatban összefoglalom a kutatási hipotézisek értékelését. Az első hipotézis igazolódott, mivel az előteszt és utóteszt közötti eredmények jelentős fejlődést mutattak. Az alkalmazást használó tanulók átlagos pontszáma 5.63 pontról 14.25 ugrott, ami jelentős fejlődést jelent. A párosított mintás t-próba eredménye szerint is ez határozottan szignifikáns volt. Így megállapíthatjuk, hogy az első hipotézisünk igazolódott.

A második hipotézis, miszerint az AI-alapú alkalmazást használó csoport fejlődése nagyobb mértékű lesz, mint a hagyományos módszerrel tanuló kontrollcsoporté nem igazolódott be. Bár a kísérleti csoport átlag fejlődése minimálisan magasabb volt, ez gyakorlatilag nem jelent szignifikáns különbséget a fejlődés mértékében. Így megállapíthatjuk, hogy nem igazolódott be a hipotézis, és a két csoport fejlődése hasonló mértékűnek tekinthető.

A harmadik hipotézis, miszerint az AI-alapú alkalmazást használó tanulók beszédkészsége javul az adott témáról folytatott irányított beszélgetés során, részben igazolódott. Az AI alapú alkalmazás lehetőséget biztosított arra, hogy a tanulók a témáról irányított beszélgetést folytassanak, aktívan használják az új szókincseket, és visszajelzésen alapuló hibajavítást kapjanak. Ez a folyamat támogatta a beszédkészség fejlődését, azonban ez a rész nem került külön statisztikai vizsgálatra, így nem lehet egyértelműen kijelenteni, hogy szignifikáns lett volna a javulás ezen a téren.

A negyedik hipotézis, miszerint a tanulók pozitívan értékelik az AI-alapú alkalmazás használhatóságát, érthetőségét és motiváló hatását, igazolódott. A visszajelzések azt mutatták, hogy az alkalmazás könnyen használhatónak bizonyult, érdekesnek tartották a tanulást új módszerrel, így pozitívnak élték meg. Ez arra utal, hogy az alkalmazás nemcsak hatékony tanulási eszköz lehet, hanem egyben motiváló és élvezetes is.

